\documentclass[11pt,oneside,leqno]{article}
\usepackage[paperwidth=6.5in,paperheight=9.5in,margin=.72in]{geometry}
\usepackage{fontspec}
\setmainfont{lmroman10-regular.otf}[BoldFont=lmroman10-bold.otf,ItalicFont=lmroman10-italic.otf,BoldItalicFont=lmroman10-bolditalic.otf,SmallCapsFont=lmromancaps10-regular.otf]
\usepackage[german]{babel}
\germansetup{glottonyms=legacy,hyphenrules=legacy}
\usepackage{amsmath,amssymb,microtype,needspace,multirow,booktabs}
\usepackage[unicode,hidelinks]{hyperref}
\hypersetup{pdftitle={Numerisches Rechnen: Anzeige und Nachlass},pdfauthor={Carl Friedrich Gauss}}
\newcommand{\SourceLeaf}[1]{\hypertarget{native-#1}{}}
\newcommand{\BareRoot}[1]{\mathord{\surd}#1}
\setlength{\parindent}{1.2em}
\setlength{\emergencystretch}{1.4em}
\widowpenalty=10000\clubpenalty=10000\raggedbottom
\begin{document}
\SourceLeaf{126}
\thispagestyle{empty}
\vspace*{.22\textheight}
\begin{center}
{\Large NUMERISCHES RECHNEN.}\par\bigskip
{\large NACHTRÄGE ZU BAND III.}
\end{center}
\clearpage
%
\SourceLeaf{128}
\begin{center}
\textbf{ANZEIGE.}
\end{center}

\medskip
\noindent\rule{\linewidth}{.3pt}

\noindent Allgemeine Literatur-Zeitung vom Jahre 1808. Halle-Leipzig 1808. Nr.\,45. Februar 12. S.\,353--358.

\noindent\rule{\linewidth}{.3pt}
\medskip

Dresden, in der Waltherschen Hofbuchhandlung: \textsc{Leonellis} \emph{logarithmische Supplemente, als ein Beitrag, Mängel der gewöhnlichen Logarithmentafeln zu ersetzen. Aus dem Französischen nebst einigen Zusätzen von} \textsc{Gottfried Wilhelm Leonhardi}, \emph{Souslieutenant beim kurfürstl. sächsischen Feldartilleriecorps.} 1806. 88 S. in Octav.

Diese kleine Schrift enthält zwei von einander unabhängige Abhandlungen, die indessen in so fern einen gemeinschaftlichen Zweck haben, als beide etwas zu leisten bestimmt sind, was sich mit den gewöhnlichen logarithmischen Tafeln ohne anderweitige Hülfsmittel entweder nicht so vollkommen, oder nicht so bequem erreichen lässt. Die eine soll nemlich dazu dienen, vermittelst einiger sehr geschmeidiger, hier zugleich mitgelieferter, Hülfstafeln, zu jeder gegebenen Zahl ohne zu grosse Mühe den Logarithmen auf mehrere Decimalen, als die gewöhnlichen Tafeln verstatten, zu berechnen. Die andere entwickelt die Idee einer besondern Tafel, vermittelst welcher die Logarithmen von Summen oder Differenzen zweier bloss durch ihre Logarithmen gegebenen Grössen durch eine einzige Operation mit einer Bequemlichkeit sollen bestimmt werden können, wie sie bei andern Verfahrungsarten nicht Statt findet; und zwar, des Vfs. Plane nach, gleichfalls mit einer ungewöhnlich grossen Anzahl von Decimalen (14); von dieser Tafel ist jedoch nur erst eine Probe beigefügt. Da das Numerische der Logarithmen und jede sich darauf beziehende Erleichterung jedem, der viel mit Zahlenrechnungen zu thun hat, von grosser Wichtigkeit ist: so wird es sich wohl der Mühe verlohnen, diesen Untersuchungen eine nicht
%
\SourceLeaf{129}
bloss oberflächliche Aufmerksamkeit zu widmen, um besonders den praktischen Werth der davon zu hoffenden Vortheile würdigen zu können.

Der Vf. hatte seine Schrift dem französischen Nationalinstitute zur Beurtheilung vorgelegt; der Bericht, welchen \textsc{Delambre} darüber abgestattet und welchen das Institut gebilligt hat, ist der Schrift selbst beigefügt. \textsc{Leonelli} ist damit nicht ganz zufrieden gewesen, und äussert in seinen gleichfalls angehängten Gegenbemerkungen seine Empfindlichkeit einigemal mit vieler Lebhaftigkeit. Wir werden nicht umhin können, auch \textsc{Delambres} Bericht und \textsc{Leonellis} Beantwortung mit zu berühren, da beide nun einen Theil des Buchs selbst ausmachen.

Bei der in der \emph{ersten} Abtheilung vorgetragenen Methode, den Logarithmen jeder vorgegebenen Grösse $A$ zu berechnen, liegt die Hauptidee zum Grunde, dass diese Grösse in ein Product von der Form
\[
10^{\mu} a \left(1+\frac{1}{10} b\right)\left(1+\frac{1}{100} c\right)\left(1+\frac{1}{1000} d\right)\ldots
\]
verwandelt werde, so dass $a$, $b$, $c$, $d\ldots$ einfache ganze Zahlen bedeuten. Man sieht erstlich, dass die Logarithmen dieser einzelnen Factoren eine abnehmende Reihe bilden werden, von denen man sich mit einer gewissen grössern oder kleinern Zahl begnügen kann, je nachdem man den Logarithmen des Products mit mehr oder weniger Decimalstellen verlangt; zweitens dass man alle jene einzelnen Logarithmen sogleich in einer ein für allemal berechneten Tafel vorräthig haben kann. Eine solche Tafel liefert \textsc{Leonelli} hier zuerst für die Briggischen Logarithmen auf 20 Decimalstellen. Diese enthält zuerst die Logarithmen der ganzen Zahlen von 2 bis 9; sodann die Logarithmen für alle $1+\frac{1}{10} b$, d.\,i. für 1,1; 1,2; 1,3 bis 1,9; nachher für alle $1+\frac{1}{100} c$ u.\,s.\,w. Ausser den 8 ganzen Zahlen kommen also, für jede Ordnung der gebrochenen, 9 Logarithmen; weiter als bis zur 11$^{\mathrm{ten}}$ Ordnung brauchte aber die Tafel nicht ausgedehnt zu werden, da bei der zwölften die Logarithmen mit den Logarithmen der eilften einerlei bedeutende Ziffern haben, und so bei den folgenden. Auf diese Weise umfasst die Tafel nur 107 Logarithmen, welche hinreichen, um den Logarithmen jedes Products der obigen Form auf 18 oder 19 Ziffern durch blosse Addition zu berechnen; die 20$^{\mathrm{ste}}$ Stelle bleibt dabei offenbar immer schwankend. Für die hyperbolischen Logarithmen hat \textsc{Leonelli} eine ähnliche Tafel von demselben Umfange beigefügt.

%
\SourceLeaf{130}
Wie sich übrigens jede Zahl $A$ unter obige Form bringen lasse, wird man leicht übersehen; $\mu$ und $a$ ergeben sich sogleich von selbst; macht man dann $A = 10^{\mu} a B$: so ist $b$ die Ziffer in der ersten Decimalstelle von $B$; setzt man ferner $B = \left(1+\frac{1}{10} b\right) C$, so ist $c$ die Ziffer in der zweiten Decimalstelle von $C$ u.\,s.\,w. Bei näherer Untersuchung findet sich, dass man noch bequemer $\frac{1}{A}$, ohne diesen Quotienten wirklich zu berechnen, unter die obige Form bringen kann, indem man successive $10^{\mu} a A = B$, $\left(1+\frac{1}{10} b\right) B = C$ u.\,s.\,w. macht, wo man also $B$ aus $A$, $C$ aus $B$ u.\,s.\,w. durch Multiplication erhält, indem bei der ersten Methode Division nöthig ist. \textsc{Leonelli} hat indessen den Calcul bei beiden Methoden durch Reduction auf einen bestimmten Mechanismus noch abgekürzt; wer häufigen Gebrauch von diesen Tafeln zu machen denkt, thut wohl, sich mit demselben vertraut zu machen. Nach einiger Übung wird man es dann gewiss bald dahin bringen, sich die anzuwendende Aufmerksamkeit mechanisch zu machen, und wir halten daher mit \textsc{Leonelli} den Vorwurf \textsc{Delambres} für übereilt, dass eine ermüdende Aufmerksamkeit nöthig sei, wenn man nicht jedesmal alle die vielen überflüssigen Nullen aufschreibe.

Man sieht leicht, dass man auf ähnliche Weise auch $A$ oder $\frac{1}{A}$ unter die Form
\[
10^{\mu} a \left(1+\frac{1}{100} b\right)\left(1+\frac{1}{10000} c\right)\ldots
\]
setzen könne, so dass $a$, $b$, $c\ldots$ lauter ganze Zahlen unter 100 bedeuten. Bei jedem einzelnen Falle wird man dann nur halb so viele Glieder nöthig haben; dagegen wird die Hülfstafel zwar auch nur halb so viel Ordnungen, aber in jeder 99 Logarithmen enthalten, also einen fast fünfmal so grossen Umfang haben müssen. \textsc{Leonelli} hat auch eine solche Hülfstafel, aber nur für die Briggischen Logarithmen, und nur auf 15 Decimalen beigefügt; sie enthält 485 Logarithmen. Es ist schade, dass nicht auch diese Tafel für 20 Decimalen eingerichtet ist.

Für die umgekehrte Aufgabe, zu einem gegebenen Logarithmen die Zahl zu finden, sind \textsc{Leonellis} Tafeln nicht weniger anwendbar. Durch Zerlegung des gegebenen Logarithmen in seine in den Tafeln befindlichen Bestandtheile erhält man die Factorenreihe der Zahl, und durch deren Multiplication (wofür sich leicht ein geschmeidiger Algorithmus findet) die Zahl selbst.

%
\SourceLeaf{131}
Wo man übrigens mit sieben oder zehn Decimalen ausreicht, gewähren freilich die gewöhnlichen oder die \textsc{Vlacq}schen, auch von \textsc{Vega} herausgegebenen Tafeln bei weitem mehr Bequemlichkeit, als die \textsc{Leonelli}schen. Und in der That sind jene auch bei den allerfeinsten astronomischen oder sonst auf die Körperwelt sich beziehenden Rechnungen, ohne Ausnahme, überflüssig genau. Bei analytischen Rechnungen, oder auch bei der primitiven Construction von Tafeln, kommt man jedoch öfters in den Fall, wo man eine grössere Schärfe wünscht, und dann, wenn man anders nicht noch mehr als 13 oder 18 Ziffern verlangt, sind \textsc{Leonellis} Tafeln unstreitig das Brauchbarste und Bequemste, was man zu diesem Behufe anwenden kann.

Dem \textsc{Delambre}schen Berichte zufolge hat übrigens \textsc{Brigg} ganz dieselbe Methode in seiner \emph{arithmetica logarithmica} vorgetragen, und eine ähnliche Hülfstafel auf 14 Decimalen berechnet; allein die französischen Geometer selbst haben von diesem seltenen Werke nur Ein Exemplar auftreiben können, worin gerade dieser Theil gefehlt hat. Das Verdienst der eigenen Erfindung und eigenen Berechnung bleibt also \textsc{Leonelli} immer ungeschmälert, und es ist billig, dass wir ihm besonders für letztere den gebührenden Dank zollen.

Den Zweck der zweiten Abhandlung haben wir bereits oben angezeigt; wir wollen nun sehen, wie \textsc{Leonelli} denselben erreichen will. Er schlägt eine aus drei Columnen bestehende Tafel vor; um uns kürzer zu fassen, wollen wir, anstatt \textsc{Leonellis} schwerfällige und unnöthige Terminologie zu gebrauchen, die zusammengehörigen Glieder dieser drei Columnen durch $P$, $Q$, $R$ bezeichnen. Diese Grössen sollen so von einander abhängen, dass, wenn man die Zahl, deren Logarithm $P$ ist, durch $x$ bezeichnet,
\[
Q = \log\left(1+\frac{1}{x}\right), \qquad R = \log(1+x)
\]
sei; daher immer $R = P+Q$ sein wird; ferner soll diese Tafel nicht nach $x$, sondern nach $P$ geordnet sein, oder $P$ soll gleichförmig wachsen, und zwar von Null bis ins Unendliche, oder vielmehr bis $Q$ als verschwindend betrachtet werden kann. Der Gebrauch einer solchen Tafel lässt sich leicht übersehen. Soll aus $\log a$ und $\log b$ der Logarithm von $a+b$ bestimmt werden: so geht man (wenn man voraussetzt, dass $a$ grösser ist als $b$) mit $\log a - \log b$ in die erste Columne ein, oder setzt diese Differenz $= P$: dann ist offenbar
\[
\log(a+b) = Q + \log a = R + \log b.
\]

%
\SourceLeaf{132}
Soll man hingegen den Logarithmen der Differenz bestimmen: so wird man $\log a - \log b$ entweder in der zweiten oder dritten Columne finden, je nachdem diese Differenz kleiner oder grösser ist, als $\log 2$. Im ersten Falle, wenn man $\log a - \log b = Q$ macht, wird $\log(a-b) = \log b - P = \log a - R$ sein; im zweiten hingegen, wenn man $\log a - \log b = \log R$ setzt, wird der gesuchte Logarithm
\[
= \log b + P = \log a - Q.
\]
Sonderbar ist’s, dass \textsc{Leonelli} diese Art, den Logarithmen der Differenz zu bestimmen, wenn die Differenz der Logarithmen kleiner als $\log 2$ ist, nicht gleich bemerkt hat; in der Schrift selbst gibt er für diesen Fall ein anderes verwickelteres Verfahren, und erst durch die Rüge dieser Unvollkommenheit in dem \textsc{Delambre}schen Berichte ist er auf die doch so nahe liegende Art, die Tafel zu benutzen, geführt, welche er in seinen Bemerkungen über diesen Bericht mit vieler Weitschweifigkeit erklärt.

Obgleich wir \textsc{Leonellis} Gedanken, durch eine solche Tafel die logarithmischen Rechnungen zu erleichtern, im Ganzen genommen unsern Beifall nicht versagen können, sondern vielmehr die wirkliche Ausführung einer solchen Tafel für wünschenswerth halten: so können wir doch allem übrigen, was \textsc{Leonelli} über diesen Gegenstand sagt, nur wenig Werth beilegen. Seine Entwickelung des Gebrauchs ist für einen so elementarischen Gegenstand mit unnöthiger Weitläuftigkeit vorgetragen. Auch nur etwas tiefere Untersuchungen über das Gesetz des Fortganges der Tafeln und der Differenzen, welche doch zur Bestimmung des ihnen zu gebenden Umfanges sehr nöthig wären, findet man gar nicht, wohl aber einige bloss hingeworfene, und zum Theil ziemlich verworren ausgedrückte Äusserungen, aus denen sich schliessen lässt, dass \textsc{Leonelli} dergleichen gar nicht, oder doch ganz unrichtig angestellt hat. Dahin gehört z.\,B. die grundfalsche Behauptung, dass, wenn $P$ immer um die Differenz 0,0005 zunimmt und sich dem Werth 6 nähert, die ersten Differenzen von $Q$ nur noch eine oder zwei bedeutende Ziffern haben sollen, wenn $Q$ mit 14 Decimalen ausgedrückt wird. Eine leichte Rechnung zeigt, dass diese Differenzen bis dahin wenigstens fünf bedeutende Ziffern behalten; noch weniger verschwinden sie, wenn $P$ den Werth 6 übersteigt, sondern erst für $P = 10,7$ werden sie bis auf eine Einheit in der vierzehnten Decimale abge\SourceLeaf{133}%
%
nommen haben; bis dahin hätte aber die Tafel über 21000 Glieder. Was aber noch wichtiger ist: man darf keinesweges die dritten Differenzen überall als verschwindend betrachten, wie \textsc{Leonelli} sich einbildet, der nur auf die zweiten Rücksicht zu nehmen für hinreichend hält. Zu diesem Irrthume scheint ihn die von ihm berechnete Probe des Anfangs der Tafel verleitet zu haben, wo freilich die dritten Differenzen verschwinden; allein eine leichte Rechnung zeigt, dass diese dritten Differenzen weiterhin \emph{zunehmen} und allerdings bedeutend werden können, wenn man 14 Decimalen geben, und keine kleinern Differenzen bei $P$, als 0,0005 gebrauchen will: für $x = 2+\BareRoot{3}$, oder für $P = 0,5719$, wo die dritten Differenzen ihr Maximum erreichen, finden wir unter obigen Voraussetzungen ihren Werth $= 6377$ Einheiten in der vierzehnten Decimale. Weit entfernt also, dass \textsc{Leonelli} mit weniger als 5000 Gliedern ausreichen könnte, müsste die Tafel, wenn die zweiten Differenzen überall hinlänglich sein sollen, eine so grosse Ausdehnung erhalten, dass ihre Berechnung die Mühe keinesweges belohnen würde. Allein wozu auch vierzehn Decimalen? Rechnungen, wo eine solche Schärfe nöthig wäre, kommen ja nur höchst selten vor, und für einen so seltenen Fall eine doch nicht sehr bedeutende Abkürzung der Arbeit zu erhalten, daran ist wenig gelegen. Hingegen für Rechnungen, die täglich vorkommen, und wo sieben Decimalen völlig hinreichen, würde eine, wenn auch nur mässige, aber oft wiederkommende Erleichterung der Arbeit allerdings schätzbar sein. Dann müsste aber die Tafel, um auch die nöthige Bequemlichkeit zu gewähren, so eingerichtet werden, dass überall auch die zweiten Differenzen verschwinden; mit etwas mehr als 15000 Gliedern liesse sich diess bequem erreichen, die nur ein mässiges Bändchen machen würden. Doch zu einer weitern Ausführung ist hier nicht der Ort.

Noch eine Probe, wie oberflächlich \textsc{Leonelli} seinen Gegenstand behandelt hat, gibt die S.\,60 vorgetragene Formel, für einen vorgegebenen Werth von $R$, der zwischen zwei der Tafel fällt, den entsprechenden von $Q$ mit Rücksicht auf die zweiten Differenzen zu finden. Diese Formel, deren Deduction dem Übersetzer, seinem Geständnisse zufolge, so viele vergebliche Mühe gemacht hat, ist ganz falsch; statt
\[
2\delta(r'n - n) \quad \text{sollte nemlich stehen} \quad \frac{(D+d)\delta rn}{d}.
\]
%
\Needspace{17\baselineskip}
\SourceLeaf{134}
\noindent Übrigens bemerken wir noch, dass, wenn man die Zahlen der ersten Columne als die doppelten Logarithmen von Tangenten betrachtet, die Zahlen der zweiten und dritten die doppelten Complemente der Logarithmen der dazu gehörigen Sinus und Cosinus sein werden. Man kann daher mit den gewöhnlichen trigonometrischen Tafeln den vorgesetzten Zweck ganz auf dieselbe Art erreichen, als mit den von \textsc{Leonelli} vorgeschlagenen; nur erspart theils die letztere die Division und Multiplication mit 2 (die aber doch ein mässig geübter Rechner leicht im Kopfe macht), theils gibt sie, wenn sie auf dieselbe Zahl von Decimalen berechnet ist, doppelt so viele Schärfe. Sonderbar ist es, dass die sich hierauf gründende Art, die trigonometrischen Tafeln zu gleicher Absicht anzuwenden, von \textsc{Delambre} nicht erwähnt wird, da er doch eine andere weit weniger bequeme anführt. \textsc{Leonelli} hat sehr Recht, sich zu beschweren, dass man ein solches Verfahren dem seinigen an die Seite setzen wollte.\par\medskip\begin{center}\rule{.24\linewidth}{.3pt}\end{center}

%

\clearpage\SourceLeaf{135}




\begin{center}
\textbf{NACHLASS.}
\end{center}

\begin{center}
[I.]
\end{center}

\begin{center}
\textbf{VORSCHRIFTEN, UM DEN LOGARITHMEN DES SINUS EINES
KLEINEN BOGENS ZU FINDEN.}
\end{center}

\[
\log\mathrm{hyp}\sin\varphi = \log\varphi - \tfrac{1}{6}\varphi\varphi - \tfrac{1}{180}\varphi^4 - \tfrac{1}{2835}\varphi^6 - \tfrac{1}{37800}\varphi^8 - \ldots
\]
\[
\log\mathrm{brigg}\sin n'' = 4{,}6855748.668 + \log n'' - An''n'';
\]
\[
\log A\ \mathrm{proxime} = -\left\{\!\begin{smallmatrix}11\\4\end{smallmatrix}\!\right\}{,}7692172 + \tfrac{1}{5}n''n''A.
\]

Noch genauer:

\[
\log A' = -\left\{\!\begin{smallmatrix}11\\4\end{smallmatrix}\!\right\}{,}7692172;
\]
\[
\log A = -\left\{\!\begin{smallmatrix}11\\4\end{smallmatrix}\!\right\}{,}7692172 + \tfrac{1}{5}n''n''\left\{A+\tfrac{17}{42}(A-A')\right\}.
\]

Hiernach findet man den Log. bis über $30^0$ auf die letzte Ziffer [7. Decimale] genau.

\[
\log\mathrm{hyp}\cos\varphi = -\tfrac{1}{2}\varphi\varphi - \tfrac{1}{12}\varphi^4 - \tfrac{1}{45}\varphi^6 - \ldots
\]
\[
\log\mathrm{brigg}\cos n'' = -Bn''n'';
\]
\[
\log B\ \mathrm{proxime}\ = -\left\{\!\begin{smallmatrix}11\\4\end{smallmatrix}\!\right\}{,}2920960 + \tfrac{1}{3}n''n''B
\]
\[
= -\left\{\!\begin{smallmatrix}11\\4\end{smallmatrix}\!\right\}{,}2920960 + \tfrac{1}{3}n''n''\left\{B+\tfrac{1}{10}(B-B')\right\}.
\]

\[
\log\mathrm{hyp}\,\mathrm{tang}\,\varphi = \log\varphi + \tfrac{1}{3}\varphi\varphi + \tfrac{7}{90}\varphi^4 + \tfrac{62}{2835}\varphi^6 + \ldots
\]
\[
\log\mathrm{brigg}\,\mathrm{tang}\,n'' = 4{,}6855748.668 + \log n'' + Cn''n'';
\]
\[
\log C\ \mathrm{proxime} = -\left\{\!\begin{smallmatrix}11\\4\end{smallmatrix}\!\right\}{,}4681872 + \tfrac{7}{10}n''n''\left\{C-\tfrac{599}{2058}(C-C')\right\}. 
\]
%
\clearpage\SourceLeaf{136}
\begin{center}
[II.]

\medskip
\textbf{[INTERPOLATION DER COTANGENTEN UND COSECANTEN KLEINER BÖGEN.]}
\end{center}

Zur Interpolation der Cotangenten und Cosecanten kleiner Bögen sind
folgende Formeln die brauchbarsten:

Es gehören
\begin{center}
\begin{tabular}{cc}
zu den Bögen & die Cotangenten (oder Cosecanten) \\
$a$ & $\theta$ \\
$x$ & $y$ \\
$a'$ & $\theta'$, \\
\end{tabular}
\end{center}
so ist
\[
\begin{aligned}
\mathrm{I.}\quad x &= a+\frac{(\theta-y)a(a'-a)}{(\theta-y)a+(y-\theta')a'}
=a+\frac{(\theta-y)(a'-a)}{\theta-\theta'+\dfrac{a'-a}{a}(y-\theta')}\\[5pt]
&=a'-\frac{(y-\theta')a'(a'-a)}{(\theta-y)a+(y-\theta')a'}
=a'-\frac{(y-\theta')(a'-a)}{\theta-\theta'+\dfrac{a'-a}{a'}(y-\theta)}.
\end{aligned}
\]
\[
\begin{aligned}
\mathrm{II.}\quad y&=\frac{a'\theta'-a\theta}{a'-a}+\frac{(\theta-\theta')aa'}{(a'-a)x}
=\theta-\frac{(\theta-\theta')(x-a)a'}{(a'-a)x}\\[5pt]
&=\theta'+\frac{(\theta-\theta')(a'-x)a}{(a'-a)x}\,.
\end{aligned}
\]

Beispiel zu II.\quad Man sucht $\mathrm{cotg}\,0{,}060335$.
\begin{center}\small
\begin{tabular}{r@{\;}r@{\quad}l@{\;}l}
$\theta=$ & $10{,}5259499$ &&\\
$\theta'=$ & $10{,}5084181$ &&\\\cline{2-2}
$\theta-\theta'=$ & $175318$ & $\log\; .\ .\ .\ .\ .\ .\ .$ & $8{,}2438265-10$\\
$x-a=$ & $0{,}000035$ & $\log\; .\ .\ .\ .\ .\ .\ .$ & $5{,}5440680-10$\\
\multicolumn{2}{r}{$[a'=0{,}0604]$} & $\log\; .\ .\ .\ .\ .\ .\ .$ & $8{,}7810369-10$\\
&& $\mathrm{Compl.}\;\log x(a'-a)\; .\ .$ & $5{,}2194307$\\\cline{3-4}
\multicolumn{4}{r}{$0{,}00614274\; .\ .\ .\ .\ .\ .\ .\ .\ .\ .\ .\ .\ .\ .\ .\ .\ .\ .\ .\ .\;7{,}7883621-10$}\\
&$10{,}5259499$&&\\\cline{2-2}
&$10{,}5198072$&&
\end{tabular}
\end{center}
%
\clearpage\SourceLeaf{137}
\begin{center}
[III.]

\medskip
\textbf{MUSTERRECHNUNG, UM AUS $A=p\cos P$, $B=p\sin P$\\
$p$ UND $P$ ZU FINDEN.}
\end{center}

\begin{center}\small
\begin{tabular}{l@{\;}l@{\qquad\hspace{.8em}}l}
$B\; .\ .\ .\ .\ .\ .$ & $9{,}56905\;69225\,n$ &\\
$A\; .\ .\ .\ .\ .\ .$ & $0{,}28523\;04177\,n$ &\\\cmidrule(r{1.4em}){1-2}
$\mathrm{tg}\,P\; .\ .\ .\ .\ .$ & $9{,}28382\;65048$ & $P=190^0\,52'\,52''{,}913567$\\
$\mathrm{tg}\,P^*\; .\ .\ .\ .$ & $9{,}28379\;34121$ & $P^*=190\;52\;50$\\\cmidrule(r{1.4em}){1-2}
$\cos P^*\; .\ .\ .$ & $9{,}99212\;15841\,n$ &\\
& $-11792$ & $5{,}5197322\; .\ .\ .\ .\ .\;3\;30927$\\
$\cos P\; .\ .\ .\ .$ & $9{,}99212\;04049\,n$ & $\underline{9{,}9842419}\; .\ .\ .\ .\;\cos P\cos P^*$\\\cmidrule(r{1.4em}){1-2}
$p\; .\ .\ .\ .\ .\ .\ .$ & $0{,}29311\;00128$ & $5{,}5039741$
\end{tabular}
\end{center}

\begin{center}\small
\renewcommand{\arraystretch}{1.35}
\begin{tabular}{r@{\;}l@{\;}l}
\multirow{2}{*}{$11792\Bigg\{$}&$9{,}2838100\; .\ .\ .\ .\;\BareRoot{(\mathrm{tg}\,P\mathrm{tg}\,P^*)}$&\\
&$\overline{4{,}7877841}$&\multirow{2}{*}{$\Bigg\}\;2''{,}913567.$}\\
$\left[-\log\frac{\mathrm{Mod.}}{206264{,}8}\right]=\mathrm{Const.}\; .\ .\ .$&$-4{,}3233592$&
\end{tabular}
\end{center}

\medskip
\begin{center}\textbf{BEMERKUNGEN.}\end{center}

Dass die Anzeige, S. 121--127, von \textsc{Gauss} herrührt, geht aus einem Briefe von
\textsc{Gauss} an \textsc{Olbers} vom 3. December 1808 hervor, worin es heisst:

\begin{quotation}\itshape
„So wünschte ich z. B. sehr, dass eine solche Tafel, wie \textup{\textsc{Leonelli}}
vorgeschlagen hat (meine Anzeige: Hallische L.Z. 1808 vom 12. Febr.), ausgeführt würde.
Für bloss 5 Decimalen habe ich selbst einmal einen Anfang gemacht. Bei solchen
Rechnungen, wo sehr viele Logarithmen von Summen oder Differenzen gesucht
werden (wie z. B. bei meiner Methode die Störungen zu berechnen), würde eine
solche Tafel eine bedeutende Erleichterung geben«.
\end{quotation}

Die Notiz [I.] ist einem Handbuch entnommen; [II.] und [III.] finden sich auf Blättern,
die den Büchern der \textsc{Gauss}-Bibliothek: \textsc{Hobert} und \textsc{Ideler},
Trigonometrische Tafeln, und \textsc{Vega}, Thesaurus Logarithmorum, angeheftet sind.
\hfill \textsc{Börsch, Krüger}.
%
\end{document}
